\chapter{Resiliensi dan Disaster Recovery}

\section*{Tujuan Pembelajaran}
\begin{itemize}
  \item Mengidentifikasi mode kegagalan dalam sistem terdistribusi dan pola resiliensi untuk menanganinya.
  \item Menetapkan RTO dan RPO berdasarkan dampak bisnis.
  \item Memilih strategi DR yang seimbang antara biaya dan risiko.
\end{itemize}

\section{Pendahuluan}

% TODO
% Cakupan: insiden PDNS 2 Surabaya pada 2024 melumpuhkan layanan ratusan instansi dan memperlihatkan bahwa cadangan yang tidak teruji sama dengan tidak punya cadangan; rumah sakit tidak boleh berhenti melayani hanya karena satu pusat data bermasalah.

\section{Mode Kegagalan dalam Sistem Terdistribusi}

% TODO
% Cakupan: kegagalan parsial, latensi tinggi, dan kegagalan berantai; titik kegagalan tunggal; kegagalan yang berasal dari penyedia atau pihak ketiga (risiko pihak ketiga dibahas rinci di IT40813).
% Kasus Medikara: gangguan layanan BPJS Kesehatan yang merambat ke antrean pendaftaran.

\section{Pola Resiliensi}

% TODO
% Cakupan: timeout, retry dengan backoff, circuit breaker, bulkhead, dan graceful degradation; contoh perilaku yang terasa oleh pengguna untuk tiap pola; risiko retry yang memperparah beban.
% Pendalaman teknis: mensimulasikan kegagalan satu layanan dan mengamati perilaku circuit breaker dengan pustaka resiliensi.

\section{Multi-AZ dan Multi-Region}

% TODO
% Cakupan: perlindungan yang diberikan multi-AZ dan multi-region; biaya dan kerumitan replikasi data; kendala lokasi data di Indonesia terhadap pilihan region cadangan (kaitan dengan Bab 2).

\section{RTO, RPO, dan Strategi Disaster Recovery}

% TODO
% Cakupan: RTO dan RPO sebagai keputusan bisnis; analisis dampak bisnis sederhana; strategi backup and restore, pilot light, warm standby, dan multi-site active-active beserta perkiraan biayanya; kelangsungan bisnis dan ISO 22301 disinggung singkat (dibahas rinci di IT40813).
% Kasus Medikara: RTO dan RPO untuk rekam medis elektronik, Medikara+, dan ERP keuangan.

\section{Ransomware dan Cadangan yang Tidak Dapat Diubah}

% TODO
% Cakupan: cara ransomware menyerang cadangan; prinsip cadangan terpisah dan tidak dapat diubah; pelajaran dari insiden PDNS 2 Surabaya 2024 tentang cakupan cadangan dan uji pemulihan.

\section{Chaos Engineering dan Game Day}

% TODO
% Cakupan: menguji resiliensi dengan sengaja memicu kegagalan secara terkendali; game day sebagai latihan bersama tim teknis dan bisnis; syarat aman sebelum memulai.

\section{Aktivitas: Rencana Resiliensi dan DR}

% TODO
% Tujuan aktivitas: menetapkan RTO, RPO, dan strategi DR untuk layanan kritis, memperkirakan biayanya, dan menyusun satu skenario uji pemulihan.
% Peran dalam tim: anggota non-teknis memimpin analisis dampak bisnis (kerugian per jam, keselamatan pasien, kewajiban regulasi); anggota teknis memetakan pola resiliensi dan arsitektur DR; seluruh anggota menyusun skenario uji.
% Langkah: (1) ambil layanan kritis dari Bab 11; (2) lakukan analisis dampak bisnis sederhana; (3) tetapkan RTO dan RPO; (4) pilih strategi DR dan perkirakan biayanya; (5) susun satu skenario uji, misalnya region utama tidak tersedia selama empat jam, beserta jadwal ujinya.
% Luaran dan tonggak capstone: C4, rencana resiliensi dan DR beserta jadwal uji pemulihan.
% Berkas sources yang direncanakan: sources/chapter12/templat-analisis-dampak-bisnis.xlsx, sources/chapter12/templat-rencana-dr.md, sources/chapter12/skenario-game-day.md

\section{Ringkasan}

% TODO

\section{Latihan}

\begin{enumerate}
  \item Jelaskan perbedaan RTO dan RPO dengan contoh layanan rekam medis elektronik.
  \item Pada insiden PDNS 2 tahun 2024, hanya sebagian kecil data yang memiliki cadangan. Pelajaran arsitektur dan tata kelola apa yang dapat diambil?
  \item Mengapa strategi \textit{multi-site active-active} belum tentu tepat untuk semua layanan Medikara Health?
\end{enumerate}
